\documentclass[12pt,a4paper,oneside]{article}
\usepackage[brazil]{babel}
\usepackage[utf8]{inputenc}
\usepackage{olymp}

\setcounter{problem}{2}

\begin{document}

\begin{problem}{Lata}{lata}{2}

A Funarbe está interessada em oferecer novos tamanhos de latas do doce de leite Viçosa.  
Faça um programa que, dadas as dimensões de uma lata cilíndrica, calcule seu volume.

%\Notes
%
%Observações, se tiver.

\InputFile

A entrada terá apenas um caso de teste.

O caso de teste é composto por dois valores reais, $R$ e $A$, respectivamente o raio e a altura da lata em centímetros. Restrições: $0 < R,A \leq 100$.

\Notes
Neste problema, considere $\pi = 3.1415$.


\OutputFile

Seu programa deve gerar apenas uma linha de saída, contendo o volume da lata em $cm^3$, com precisão de 2 casas decimais.

Não se esqueça de quebrar a linha após a impressão do valor.

\Example

\begin{example}
\exmp{
5 20
}{
1570.75
}%
\end{example}

\begin{example}
\exmp{
3 25
}{
706.84
}%
\end{example}

\begin{example}
\exmp{
7 45
}{
6927.01
}%
\end{example}


%Comentários sobre o exemplo, se necessário.

\end{problem}

\end{document}
